% One Britannica scan as extracted by Nemotron Parse 3.5. Filled by src/nemotron_challenge/stages/latex.py.
% The TITLE, BODY and SOURCE markers (each wrapped in two percent signs) are replaced. Compile with XeTeX (tectonic).
\documentclass[10pt]{article}
\usepackage[a4paper,margin=18mm]{geometry}
\usepackage{fontspec}
\usepackage{multicol}
\usepackage{graphicx}
\usepackage[hidelinks]{hyperref}
\usepackage[normalem]{ulem}
\setlength{\parindent}{0pt}
\setlength{\parskip}{4pt}
\setlength{\columnsep}{8mm}
\begin{document}
\section*{Historical source: FHOU}
\begin{multicols}{2}
\subsection*{\textit{continued:} THOTHR}
Thoos. attended, walking as mourners. The queen and princesses also assisted at the \sout{ceremony} \textbf{ceremony, and the concourse} of \sout{people} \textbf{people,} including the public bodies, \sout{mountains} \textbf{municipal} and \sout{ascidental, amoured} \textbf{academical, amounted} to many thousands. An interesting part of the ceremony was the performance of a \sout{woman} \textbf{requiem} written by the \sout{sculptor’s initiale} \textbf{sculptor's intimate} and \sout{the} \textbf{dear} friend \sout{of Manchester. Thowments} \textbf{Ehlenschlager. Thorwaldsen} was rather above the middle height, firmly \sout{build,} \textbf{built,} with broad shoulders. The \sout{owner} \textbf{outline} of his face was rather square. His general expression was calm and thoughtful, and very pleasing. His eyes were light and penetrating; his mouth wide and usually \sout{enjoyable.} \textbf{closely shut.} He wore his hair, \sout{and} \textbf{which} had grown nearly \sout{white} \textbf{white,} in large masses divided over his fine \sout{than forefoot,} \textbf{broad forehead,} and \sout{killing} \textbf{falling} on each side, something like a \sout{lion’s} \textbf{lion's} mane, nearly to his shoulders. Thornflies we all always ask me walk van Cuavaa and Flacman as one of the great representatives of Sculpture. The former, his sliders were almost the first to break through the intensive routine of the world mixed Italian and French art which had descended from Bernent and his school. Thornalizes nobly asked in advancing this revoltion, and with them assisted in establishing the practice of sculpture upon true principles. (a. w.-r.).

\subsection*{FHOU}
{\small\itshape [Margin: 239]}\par
\sout{FHOU, Agorge Actors Drz,} \textbf{THOU, JACQUES AUGUSTE DE,} or, as \sout{he is} \textbf{his} name is \sout{for quantity} \textbf{frequently} written in its \sout{Ladurized} \textbf{Latinized} form, \sout{Acores Atorures Trnucts,} \textbf{JACOBUS AUGUSTUS THUANUS,} was \textbf{a} distinguished French \sout{historians,} \textbf{historian,} and was born at Paris on the \sout{860} \textbf{8th} of October 1553. In his infancy he was very weak and sickly, and \sout{infinity} \textbf{infirmity} of health \sout{dung} \textbf{clung} to him into \sout{many} \textbf{man's} years. \sout{Indigenous} \textbf{Indisposed} to play like other boys by reason of his \sout{final} \textbf{frail} health, \sout{be} \textbf{he} cultivated a taste \sout{of} \textbf{for} drawing and reading when \sout{she} \textbf{he} should have been engaged in the boyish \sout{amazonas} \textbf{amusements} peculiar to his years. \sout{Do Two were} \textbf{De Thou was} originally designed for the church, but he allowed his mind to take a much wider \sout{experience} \textbf{course} in \sout{his third} \textbf{its thirst} for knowledge than was fairly consistent with an aspirant after \sout{private} \textbf{priest's} orders. \sout{De greatly} \textbf{He greedily} studied literature and science, and contracted so great \sout{in} \textbf{a} liking \sout{it on Gujia,} \textbf{for Cujas,} the celebrated \sout{junk,} \textbf{jurist,} that he set out for \sout{VANCE} \textbf{Valence} in \sout{Dupphins,} \textbf{Dauphiné,} and attended his lectures on Pajinian. Here he formed the \sout{acquainness} \textbf{acquaintance} of Joseph \sout{Geliery,} \textbf{Scaliger,} who remained one of his constant friends \sout{throughout-out} \textbf{throughout} life. De \sout{Thus} \textbf{Thou} returned to Paris in 1572 after a year's absence, to witness the marriage of Henry of \sout{Warwers} \textbf{Navarre;} to see the \sout{ideal} \textbf{dead} body of \sout{College} \textbf{Coligny} hanging from the \sout{gribed;} \textbf{gibbet;} and to be a spectator of the horrible massacre of St \sout{Bartholne louner.} \textbf{Bartholomew.} Next year he went to Italy as an \sout{attention} \textbf{attendant} of Paul de Foix; \sout{her} \textbf{he} visited \sout{Panderand} \textbf{Flanders} and \sout{Elanall} \textbf{Holland} in 1576, and in 1578 he accepted, with reluctance, the office of an ecclesiastical \sout{condition} \textbf{councillor} of the \sout{patience} \textbf{parlement} of Paris. In 1582, while on a visit to Bordeaux, he formed the \sout{friend-ship} \textbf{friendship} of honest, \sout{garnished} \textbf{garrulous} old \sout{Monatarie,} \textbf{Montaigne,} whose character and genius he has warmly praised. His \sout{Harry} \textbf{father,} Christophe de Thou, first president of the \sout{patented} \textbf{parlement} of Paris, having died during the same year, and both his elder and younger brothers being now dead, it was deemed advisable for \sout{the students of her parents. One of how} \textbf{De Thou} to cease to be \sout{recalculated.} \textbf{an ecclesiastic.} He \sout{recently assigned} \textbf{accordingly resigned} that \sout{Professors} \textbf{profession} soon after, and accepted \sout{their effect} \textbf{the office of} Master of Requests, and was appointed to his \sout{uncle’s} \textbf{uncle's} place of \sout{pri- \_ident do matter} \textbf{président à mortier} to the \sout{patient} \textbf{parlement} of Paris. The \sout{fact Thoms} \textbf{De Thous} had always remained firm \sout{adherence} \textbf{adherents} of \sout{reality,} \textbf{royalty,} and the most \sout{illustriations} \textbf{illustrious} of them would not \sout{disclose the interest.} \textbf{disgrace his ancestry.} He was admitted a \sout{counselled or} \textbf{counsellor of} state to Henry \sout{II:} \textbf{III.} at Chartres, and was engaged in various delicate and \sout{importants} \textbf{important} missions for Henry \sout{III:} \textbf{III.} during the succeeding years. He visited \sout{Bolsa,} \textbf{Blois,} Paris, Germany, and \sout{later,} \textbf{Italy,} all in the royal cause. The death of Henry reached him in Venice in \sout{1584.} \textbf{1589.} De Thou immediately \sout{sat} \textbf{set} out \sout{in} \textbf{for} France by way \textbf{of} Switzerland, and was received very graciously by Henry \sout{IV:} \textbf{IV.} at \sout{Chakaranlus. Henecourtard} \textbf{Château-dun. Henceforward} he became a very useful servant to the crown. He was made keeper of the Royal Library on the death of Bishop \sout{Amovit} \textbf{Amyot} in 1591, and in 1593 he commenced his great work Historia \sout{Sul} \textbf{Sui} Temporis, on which he had meditated for fifteen years. From this period until his death his time was much \sout{rec-ppled} \textbf{occupied} in the composition of his history, and except the arrangement of the \sout{Zuez} \textbf{Edict} of \sout{Nations,} \textbf{Nantes,} in which he had an important share, he engaged very little in public business. The death of Henry in 1610 \sout{eventually} \textbf{virtually} brought his political life to \textbf{a} termination. \sout{Seevent} \textbf{Sourced} with \sout{high disappointment,} \textbf{public disappointments,} and \sout{ten} \textbf{torn} by private \sout{greids,} \textbf{griefs,} caused mainly by the death of his second wife, he died at Paris on the 7th of May 1617, in his sixty-fourth year. A son of De \sout{Thous} \textbf{Thou's} by his second wife, François August, \sout{full} \textbf{fell} a sacrifice to the \sout{inexor-able vegences} \textbf{inexorable vengeance} of Cardinal \sout{Richardson, and to} \textbf{Richelieu on the} 12th of \sout{Septtender} \textbf{September} 1642. From 1604, when the first night books of De Than’s History appeared, the author was received with declaration by all the reading public of Europe, as the first historian of the age. During this life there appears everybody back of it in all, and the remainder, forming in all one hundred and thirty-eight books, were published in 1620. The space view which his record extends, if from one 1644 to 1607; from the closing years of Francis I. till nearly the completion of the reign of Henry IV. It has all the fril-ness and freshness of a marative, written as a man who was a principal actor in six scenes which is his business to chronicle; it has, besides, the rare month of being argueably impacted, the author recognising with uncommon framhances the excellence and ever economic superiority of racismen and citizens who were moved rists to the nation which it was both his duty and his pride to honour and Lara. But it is much more upon the excellence of the workmanship that De Than’s History lives, then upon the counterparty accuracy and fulfillness of its recorded facts. While the style of the work certainly displays abonance of easy, flowing clouence, it is perhaps delicient in the highest quality of vivid judgements. The language in which it was written was Latini and the author has pared neither time nor honor to make its severely classical. That is does not come quite up to this ideal is perhaps, not to be wondered at... The best English edition of Thous. History is that in 7 vols by Samuel Buckley, 1723. De Thion likewise works on number of Latin papers, one of which, entitled \_De Rio Accipitraria\_ (On Hawking), prepared in 1684. The Life of this author has been freqently written. The advice noted that has gathered from an autobiography, which remains an 1601, and from Collinans \_Life of Thousand\_, London 1897. \sout{THOUROUT} \textbf{THOUROUT,} a town of Belgium, capital of a canton in the province of West \sout{Banker, I} \textbf{Flanders, 11} miles \sout{By,} \textbf{S.W.} of \sout{Burgos-} \textbf{Bruges.} It has a fine church, \sout{is townhall,} \textbf{a town-hall,} several schools, and \sout{bhec-volent} \textbf{benevolent} institutions. \sout{Lion,} \textbf{Linen,} hats, leather, \sout{goodage, entharsure, spans, chances,} \textbf{cordage, earthenware, soap, tobacco,} and other goods \textbf{are} made here. In \textbf{the} vicinity are the remains of an old castle, once the \sout{restdance} \textbf{residence} of the \sout{Count} \textbf{Counts} of Flanders. Pop. \sout{8465.} \textbf{8405.}

\subsection*{THIROCE}
\sout{THIROCE,} \textbf{THRACE,} the ancient \sout{stone} \textbf{name} of a country in the S.E. of Europe, now forming part of European Turkey. In the earliest times the name was \sout{sent} \textbf{used} in a \sout{vastay seen,} \textbf{very vague sense,} and applied to the whole region \sout{such} \textbf{north} of the \sout{Afghan Sen,} \textbf{Aegean Sea,} including Macedonia and even \sout{Seythisu} \textbf{Scythia;} but \sout{with} \textbf{when} this \sout{morning become} \textbf{country became} better known to the \sout{Greeky. Three} \textbf{Greeks, Thrace} was confined within narrower and more definite limits. It was bounded on the \sout{7b,} \textbf{N.} by the Danube, E. by the \sout{Exiton} \textbf{Euxine} and the \sout{Boskolome,} \textbf{Bosphorus,} S. by the \sout{Proponents} \textbf{Propontis} and the Aegean, and W. by the Strymon, or, \sout{as} \textbf{at} a later \sout{period} \textbf{period,} by the \sout{Neustus, and the} \textbf{Nestus, which} separated \sout{in} \textbf{it} from Macedonia. After the Roman conquest, its size was still further reduced by the chain of \sout{the same} \textbf{Haemus} being made its northern boundary. It corresponds \sout{parting} \textbf{partly} with the modern \sout{Rumell,} \textbf{Rumelia,} which also includes part of Macedonia and Thessaly. (See \sout{Rutuz)} \textbf{Rumelia.)}

\subsection*{THRARESTON}
\sout{THRARESTON,} \textbf{THRAPSTON,} a \sout{market-own} \textbf{market-town} of England, in the county and \sout{13} \textbf{18} miles \sout{ANNEX} \textbf{N.N.E.} of Northampton, on the Nen, which is here crossed by a bridge. The parish church is a cruciform building with a tower and spire; and there are
\end{multicols}
\section*{Source record}
\begin{itemize}
\item Page ID: \texttt{britannica:309f765b1ecc83e0a7f7762e44b4a927a7555443:train:604}
\item Scan: \texttt{page-0604-c50a825dd732.jpg}, 600 x 807 px
\item Image SHA-256: \texttt{\small c50a825dd732a96ef099cb60b5160d5fadaf7974163160bf5a2ac4446c70b631}
\item Extraction: \texttt{nvidia/nemotron-3.5-parse}, raw reply \texttt{artifacts/extract/raw/0604.json}
\item Text is Parse OCR output, normalized (long s as s, line-end hyphens joined). It is not corrected by hand.
\item \sout{Struck-through} words are Parse's reading; \textbf{bold} words are Super VL's correction, checked against the scan image. Both are model output.
\end{itemize}
\end{document}
